% =============================================================================
% Lecture 8: Sensor Perception
% Physical AI: Machine Learning Systems That Sense and Act (Volume IV), Chapter 8
% =============================================================================
% Thesis: an observation serves a moving machine only when it states when it
% was captured, in which frame it was measured, and how far its error has been
% measured, because the machine travels through every millisecond of its age.
% Hook: Williston, Florida, 2016 (NTSB HAR-17/02). Running case: the warehouse
% mobile manipulator at its 1.5 m/s drive limit, a worker at a rack end 1.10 m
% ahead. Real data: KITTI raw logs (Velodyne HDL-64E and four cameras),
% MLPerf Inference v3.1 on Jetson AGX Orin, the Ouster OS1 datasheet, ROS 2
% message definitions.
%
% Build:   make deck V=4 D=08_perception   (figures: images/make_figs.py)
% Numbers: images/calcs.py prints every derived value on these slides.
% Plan:    lecture_map.yml in this folder.
% =============================================================================

\documentclass[aspectratio=169, 12pt]{beamer}
\usepackage{../../assets/beamerthememlsys}
\mlsyssetup{
  volume={Volume IV},
  chapter={Chapter 8},
  chaptertitle={Sensor Perception},
  layer=boundary_brain,
  logo={../../assets/img/logo-mlsysbook.png},
  instlogo={../../assets/img/logo-harvard.png},
}
\usepackage[T1]{fontenc}
\usepackage[scaled=0.9]{helvet}
\usepackage{courier}
\renewcommand{\familydefault}{\sfdefault}
\usepackage{booktabs,amsmath}
\graphicspath{{images/}}
% Generated by calcs.py. Do not edit; edit calcs.py and rerun it.
\defval{v_drive}{1.5}
\defval{d_per_ms}{1.5}
\defval{d_clear}{1.10}
\defval{d_clear_mm}{1{,}100}
\defval{a_brake}{2}
\defval{delta_loc}{50}
\defval{delta_margin}{100}
\defval{side_clearance}{150}
\defval{t_onset}{20}
\defval{yaw_rate}{0.5}
\defval{v_aisle}{1.3}
\defval{d_braking}{562.5}
\defval{d_overhead}{150}
\defval{d_onset}{30}
\defval{t_lease_path}{62}
\defval{d_lease_path}{93}
\defval{d_rows14}{835.5}
\defval{d_spare14}{264.5}
\defval{williston_mph}{74}
\defval{williston_mps}{33.1}
\defval{t_exp}{16}
\defval{d_exposure}{24}
\defval{t_readout}{16.6}
\defval{t_arrival}{26.6}
\defval{d_arrival}{39.9}
\defval{t_trigger}{8}
\defval{d_trigger}{12}
\defval{t_link_dma}{2.0}
\defval{misorient}{0.5}
\defval{kitti_sweeps}{111}
\defval{kitti_sweep_ms}{103}
\defval{kitti_fwd}{51.6}
\defval{kitti_color_off}{10.5}
\defval{kitti_gray_off}{16.7}
\defval{kitti_gap}{6.2}
\defval{kitti_gap_whole}{6}
\defval{kitti_per_lidar_min}{103.1}
\defval{kitti_per_lidar_max}{103.3}
\defval{kitti_per_color_min}{103.0}
\defval{kitti_per_color_max}{103.3}
\defval{kitti_per_gray_min}{100.4}
\defval{kitti_per_gray_max}{105.8}
\defval{hdl_vfov}{26.8}
\defval{kitti_sweep0}{103}
\defval{kitti_v}{13.2}
\defval{kitti_kmh}{47}
\defval{kitti_sweep0_travel}{1.36}
\defval{kitti_points}{121{,}015}
\defval{kitti_gray_travel_cm}{22}
\defval{rows}{3{,}040}
\defval{t_row}{5.46}
\defval{d_skew}{24.9}
\defval{d_lateral}{9.1}
\defval{lateral_pct}{6.1}
\defval{d_skew_aisle}{21.6}
\defval{f_imu}{1{,}000}
\defval{f_lidar}{20}
\defval{t_sweep}{50}
\defval{d_sweep}{75}
\defval{rot_sweep}{25}
\defval{angle_err}{1}
\defval{d_deg_rack}{19.2}
\defval{mount_range}{15}
\defval{d_deg_far}{262}
\defval{mount_over_side}{1.7}
\defval{link_len}{1}
\defval{link_dt}{10}
\defval{d_link}{0.23}
\defval{t_st_exphalf}{8}
\defval{t_st_readout}{16.6}
\defval{t_st_link}{0.8}
\defval{t_st_dma}{1.2}
\defval{t_st_isp}{2.5}
\defval{t_st_backbone}{22}
\defval{t_st_ipc}{0.5}
\defval{t_capture}{24.6}
\defval{t_through_isp}{29.1}
\defval{t_age}{51.6}
\defval{d_age}{77.4}
\defval{pct_backbone}{43}
\defval{pct_capture}{48}
\defval{rn_p99}{11.79}
\defval{rn_p99_short}{11.8}
\defval{backbone_over_rn}{1.9}
\defval{capture_over_rn}{2.1}
\defval{t_tau}{133.6}
\defval{t_rest}{82}
\defval{d_rows15}{912.9}
\defval{d_spare15}{187.1}
\defval{t_sync}{0.2}
\defval{t_dispatch_limit}{51.8}
\defval{f_frame}{60}
\defval{t_frame}{16.7}
\defval{d_frame}{25.0}
\defval{frame_pct}{13.4}
\defval{batch_words}{four}
\defval{held_frames_words}{three}
\defval{t_batch_wait}{50}
\defval{d_batch_wait}{75}
\defval{t_batched}{101.6}
\defval{pm_deg}{45}
\defval{pm_rad}{0.785}
\defval{wc}{10}
\defval{f_c}{1.6}
\defval{t_pm}{78.5}
\defval{t_pm_used}{60}
\defval{pm_used_pct}{76.4}
\defval{st_b}{0.20}
\defval{st_f}{800}
\defval{st_sd}{0.5}
\defval{st_rack}{3.8}
\defval{st_z_loc}{4.0}
\defval{st_z_side}{6.9}
\defval{st_z_side_half}{4.9}
\defval{st_range_ratio_words}{five}
\defval{st_growth}{25}
\defval{st_near}{2}
\defval{st_far}{10}
\defval{st_sigma_near}{12.5}
\defval{st_sigma_far}{312.5}
\defval{inset_near}{162.5}
\defval{inset_far}{462.5}
\defval{os1_prec_min}{5}
\defval{os1_prec_max}{30}
\defval{os1_prec_min_cm}{0.5}
\defval{os1_prec_max_cm}{3}
\defval{os1_acc}{25}
\defval{os1_acc_retro_cm}{5}
\defval{os1_ptp}{1}
\defval{os1_drift_ppm}{10}
\defval{drift_us}{10}
\defval{sync_interval}{1}
\defval{alpha_pct}{1}
\defval{n_min}{99}
\defval{n_resid}{1{,}000}
\defval{rank}{991}
\defval{rank_top}{10}
\defval{f_vlm}{2}
\defval{sweeps_per_vlm_words}{ten}
\defval{sat_pct}{7.5}
\defval{pix_max}{255}
\defval{fog_beta}{0.05}
\defval{fog_eps}{0.05}
\defval{d_fog}{60}

% Volume IV figure roles (palette.md), for text beside the figures
\definecolor{figsense}{HTML}{007A87}
\definecolor{figpropose}{HTML}{2563EB}
\definecolor{figpermit}{HTML}{059669}
\definecolor{figveto}{HTML}{DC2626}
\definecolor{figamber}{HTML}{D97706}
\setcounter{mlsystotalsections}{5}

\title{Sensor Perception}
\author{Vijay Janapa Reddi}
\institute{Harvard University}
\date{}

% A large answer choice for predict and concept-question slides
\newcommand{\choice}[2]{\textbf{\textcolor{mlsysaccent}{#1}}\enspace #2}
% A choice that is highlighted from overlay #1 on (the answer)
\newcommand{\answerat}[3]{\alt<#1->{\colorbox{mlsysaccent!15}{\choice{#2}{#3}}}{\choice{#2}{#3}}}
% A goal or takeaway line with its text-free thumbnail
\newcommand{\thumbline}[2]{%
  \begin{minipage}[c]{2.6cm}\centering\includegraphics[width=2.5cm,height=0.95cm,keepaspectratio]{#1}\end{minipage}%
  \hspace{0.3cm}\begin{minipage}[c]{0.74\textwidth}\small #2\end{minipage}}
% One line of a worked example: shown from overlay #1 on
\newcommand{\wstep}[2]{\uncover<#1->{#2}\\[0.2cm]}
% An exact excerpt of real tool output, a log, or source, at a legible size
\newenvironment{logblock}[1][8.5]{\par\ttfamily\bfseries\fontsize{#1}{#1}\selectfont\setlength{\baselineskip}{1.25\baselineskip}\raggedright\setlength{\parskip}{0pt}}{\par}
% A photo with a one-line caption under it
\newcommand{\urow}[4]{\uncover<#1->{#2} & \uncover<#1->{#3} & \uncover<#1->{#4}\\}
\newcommand{\photocap}[3]{\begin{minipage}[t]{#1}\centering\safeimg[width=\linewidth,height=2.0cm,keepaspectratio]{#2}\\[0.05cm]{\footnotesize #3}\end{minipage}}

\begin{document}

% =============================================================================
% TITLE AND HOOK
% =============================================================================
\mlsysnote{[0.5 min] [CORE] Read the thesis aloud once. Today opens Part III of the course, Running the Machine, whose law is ``proposal is not permission.'' Perception is the first runtime stage. It turns photons and laser returns into claims a moving machine acts on. The notes total about 72 minutes; the core alone is about 58. FLEX: notes tagged [OPTIONAL] can be cut; cut first the ROS header, the phase-margin slide, and the glare slide. The planned second lecture splits after \S8.4, at the Error divider: capture and age in the first, error and the observation contract in the second.}
\mlsystitle{Sensor Perception}{An observation serves a moving machine only when it states when it was captured, in which frame, and with what measured error.}{cover_perception_blueprint.png}

\begin{frame}{Predict: what failed at Williston?}
  \mlsysactivity{Predict}{3 min}
  \begin{columns}[c]
    \begin{column}{0.55\textwidth}
      \centering\safeimg[width=\textwidth,height=4.6cm,keepaspectratio]{ntsb-williston.png}
    \end{column}
    \begin{column}{0.43\textwidth}
      \small Florida, May 2016. A car on Autosteer at \textbf{\val{williston_mph} mph} meets a semitrailer turning across its path. No warning, no braking.\\[0.25cm]
      \choice{A}{the camera stopped working}\\[0.08cm]
      \choice{B}{the radar stopped working}\\[0.08cm]
      \answerat{2}{C}{all parts ran; none saw the trailer}\\[0.08cm]
      \choice{D}{the trailer was detected, too late}
    \end{column}
  \end{columns}
  \mlsyscredit{NTSB, Highway Accident Report HAR-17/02 (2017), Findings 1 and 3; board-meeting presentation, photos courtesy of Florida Highway Patrol.}
  \note{[3 min] [CORE] LINK: this is the crash whose NHTSA report opened Lecture 7. NARRATE: the photos are from the NTSB board meeting: the trailer's right side, and the car that passed beneath it. The car's speed just before impact was \val{williston_mph} mph (\val{williston_mps} m/s, derived), Traffic-Aware Cruise Control and Autosteer active. ENGAGE: vote with fingers, one minute. Step 2: C. Finding 1: no mechanical system on either vehicle failed. Finding 3: the automated system ``was not designed to, and did not, identify the truck crossing the car's path''; no alert, no automatic braking. The NHTSA investigation reported that braking required camera and radar to agree. WARN: students who pick D assume perception either sees or is late; the third case, a silent omission with every health signal green, is the hardest one and returns in the last segment. LINK: hold the second question for the rest of the lecture: even when perception does see the obstacle, is the observation still true by the time the machine acts on it? FLEX: [CORE].}
\end{frame}

\begin{frame}{A correct detection of a worker can still arrive too late}
  \begin{columns}[c]
    \begin{column}{0.40\textwidth}
      \centering\safeimg[width=\textwidth,height=4.1cm,keepaspectratio]{er-flex-mobile-manipulator.jpg}\\[0.05cm]
      {\footnotesize A real machine of the running machine's class: a wheeled base (MiR250) with an arm (UR16e).}
    \end{column}
    \begin{column}{0.58\textwidth}
      \centering\safeimg[width=\textwidth,height=4.0cm,keepaspectratio]{aisle.pdf}\\[0.1cm]
      \small The camera names the worker correctly. Where was the worker when the light left?
    \end{column}
  \end{columns}
  \mlsyscredit{Photo: ER-FLEX by Enabled Robotics, Auledas, CC BY 4.0, Wikimedia Commons. Aisle: Sensor Perception, \S8.1 and Napkin Math 8.2 (illustrative values).}
  \note{[2 min] [CORE] LINK: the running case of the course, the warehouse mobile manipulator, a wheeled base with an articulated arm, navigation lidar, and cameras. The photo shows a commercial machine of the same class; the course's machine is a reference platform with illustrative numbers. NARRATE: the base drives an aisle at its \val{v_drive} m/s drive limit. A worker steps out at a rack end \val{d_clear} m ahead. Every millisecond is \val{d_per_ms} mm. A vision model can classify the worker with near-perfect accuracy and still hand the planner a position that is out of date. ENGAGE: what three things must a reported position say before the machine may act on it? Take answers; the expected three are when it was captured, in which coordinate frame, and how large its error is. WARN: machine learning students equate a good detection with a good observation; benchmark accuracy says nothing about time. FLEX: [CORE].}
\end{frame}

\begin{frame}{Recall: what does the stopping budget already spend?}
  \mlsysactivity{Retrieval}{2 min}
  {\normalsize\bfseries Question of the day}\enspace{\small Which stopping-budget segments are named so far?}\par\vspace{0.05cm}
  {\centering\uncover<2->{\mlsyskitfig[width=0.95\textwidth,height=3.7cm,keepaspectratio]{budget_bar_step4}\par}}\vspace{-0.15cm}
  \uncover<2->{\footnotesize Braking \val{d_braking}, allowances \val{d_overhead}, onset \val{d_onset}, and lease path \val{d_lease_path} make \textbf{\textcolor{mlsysaccent}{\val{d_rows14} of \val{d_clear_mm} mm}}.}\par\vspace{0.05cm}
  \uncover<3->{\footnotesize Lecture 7 also left the runtime two perception monitors: input support, and latency and freshness.}
  \mlsyscredit{Kit bar; Napkin Math 4.1; the evaluation record (Lecture 7).}
  \note{[2 min] [CORE] Retrieval, 60 seconds of writing; ask aloud also what Lecture 7's evaluation record left to the runtime. Step 2: the kit bar as Lecture 4 left it, and the four segments named so far at \val{v_drive} m/s. Braking at the credible deceleration of \val{a_brake} m/s$^2$: \val{d_braking} mm. Localization allowance \val{delta_loc} mm plus protective clearance \val{delta_margin} mm. Brake onset, \val{t_onset} ms: \val{d_onset} mm. Lecture 4's chunk lease plus one tick and one bus cycle, \val{t_lease_path} ms: \val{d_lease_path} mm. Sum \val{d_rows14} mm; \val{d_spare14} mm of the \val{d_clear} m remains (calcs.py). Step 3: Lecture 7 closed Part II with the evaluation record, which hands the runtime a list of obligations; two of its runtime monitors, the input-support monitor and the latency and freshness monitor, can act only on observations that carry a capture time, a frame, and an uncertainty. LINK: today names one more segment inside the delay, the travel while the observation ages. FLEX: [CORE]; a class that just had Lecture 7 can skip the writing.}
\end{frame}

\begin{frame}{Why can an accurate detection still be too late to guide a robot?}
  \thumbline{thumb_guard.pdf}{Choose a sensor by the physical quantity the decision needs, and say which one may veto.}\\[0.18cm]
  \thumbline{thumb_capture_step3.pdf}{Date a measurement at capture and place it through calibrated frames.}\\[0.18cm]
  \thumbline{thumb_waterfall_step4.pdf}{Build observation age stage by stage, \val{t_age} ms here, and charge it to the stopping budget.}\\[0.18cm]
  \thumbline{thumb_inset_step2.pdf}{Set a clearance inset from measured residuals, and write the observation contract.}\\[0.1cm]
  {\footnotesize\textcolor{midgray}{Thread: the third law, proposal is not permission. Record: the observation contract.}}
  \note{[1 min] [CORE] The driving question, then four goals, one per stretch of the lecture; each thumbnail is the shape of the figure that teaches it. The bottom line names the thread. Lecture 1 named the four laws; today opens Part III and states the third, ``proposal is not permission,'' a design principle the book elevates. A confident detection is evidence behind a proposal, and the permission path can refuse it only if the evidence carries an age, a frame, and an error bound it can check. Today adds the sixth record of the chain Lecture 1 introduced, the observation contract.}
\end{frame}

% =============================================================================
% SECTION: SENSORS
% =============================================================================
\mlsysnote{[0.5 min] [CORE] First question: which physical quantity does each decision need, and which sensor measures it?}
\mlsysdivider{Which sensor measures what the decision needs?}{Sensor Perception, \S8.1 and \S8.2.3 to \S8.2.6}
\section{Sensors}

\begin{frame}{Each sensor measures one physical quantity}
  \centering
  \photocap{2.75cm}{rpi-hq-camera-imx477.jpg}{\textbf{Camera}\\ bearing, appearance}\hfill
  \photocap{2.75cm}{realsense-d435.jpg}{\textbf{Active stereo}\\ depth by triangulation}\hfill
  \photocap{2.75cm}{velodyne-hdl64e.jpg}{\textbf{Spinning lidar}\\ range from pulse return time}\hfill
  \photocap{2.75cm}{imu-itg3200-adxl345.jpg}{\textbf{IMU (inertial)}\\ rotation rate, acceleration}\hfill
  \photocap{2.75cm}{prophesee-event-camera.jpg}{\textbf{Event camera}\\ change in brightness}\\[0.35cm]
  \small A single camera pixel fixes a direction and leaves the range open.
  \mlsyscredit{Wikimedia Commons: RPi HQ camera, Laserlicht, CC BY-SA 4.0; RealSense D435, M. Auledas, CC BY-SA 4.0; HDL-64E, S. Jurvetson, CC BY 2.0; SparkFun, CC BY 2.0; Prophesee, Auledas, CC BY 4.0.}
  \note{[2.5 min] [CORE] Real devices, left to right. A camera, here the Raspberry Pi HQ module with the Sony IMX477, the sensor the chapter's readout budget cites: bearing and appearance, no range. An active stereo camera, Intel's D435: it projects infrared speckle and triangulates, and its error grows with the square of range. A spinning lidar, the 64-beam Velodyne HDL-64E: each laser pulse returns a range directly, $r = c\,t/2$. An inertial measurement unit (IMU): angular rate and specific force, the motion between frames. An event camera: each pixel reports a change in log brightness, so a worker who stands still makes no events. Tactile sensors and radar, not shown, measure contact and Doppler velocity. ENGAGE: which of these could see a worker who has stopped at the corner, in dim light, at 4 m? (Expected: lidar; stereo needs texture and light; the event camera loses a static target.) WARN: more sensors do not help unless their failures are independent. FLEX: [CORE].}
\end{frame}

\begin{frame}{The camera supplies appearance; the lidar supplies range}
  \centering
  \safeimg[width=14.6cm,height=4.6cm,keepaspectratio]{kitti_overlay.pdf}\\[0.15cm]
  \small One real Velodyne sweep, \val{kitti_points} returns, projected into the color camera.
  \mlsyscredit{KITTI raw data, drive 2011\_09\_26\_0001, frame 0, and its calibration (Geiger et al., IJRR 2013; CC BY-NC-SA 3.0). Projection: ours.}
  \note{[2 min] [OPTIONAL] Real data from KITTI, a public driving dataset recorded in Karlsruhe in 2011 with an HDL-64E and four cameras. Each dot is a lidar return, colored by its range; the image under it is the left color camera's frame. NARRATE: the camera tells you there are parked cars and a tram track; only the dots tell you that the nearest car is about 8 m away. The projection uses three calibrated transforms: lidar to camera, rectification, and the camera's intrinsic matrix; the next segment returns to that chain. ENGAGE: why do the dots stop above a certain height in the image? (Expected: the lidar's vertical field of view is \val{hdl_vfov} degrees; above it there are no returns, only pixels.) WARN: the camera image is dense and the lidar is sparse, yet only the sparse channel measures range. FLEX: [OPTIONAL]; cut if short, the previous slide carries the idea.}
\end{frame}

\begin{frame}{The camera proposes, the lidar admits, a scanner stops}
  \begin{columns}[c]
    \begin{column}{0.62\textwidth}
      \centering\safeimg[width=\textwidth,height=5.4cm,keepaspectratio]{guard.pdf}
    \end{column}
    \begin{column}{0.36\textwidth}
      \centering\safeimg[width=\textwidth,height=3.0cm,keepaspectratio]{robotino-sick-scanner.jpg}\\[0.05cm]
      {\footnotesize A planar laser scanner (yellow, SICK) on a small mobile robot.}\\[0.2cm]
      {\footnotesize ISO 3691-4: detect persons at Performance Level d.}
    \end{column}
  \end{columns}
  \mlsyscredit{Sensor Perception, \S8.1 (ISO 3691-4, ISO 13849-1, IEC 61496-3). Photo: Robotino with SICK scanner, Robodoc9, CC BY-SA 3.0, Wikimedia Commons.}
  \note{[2.5 min] [CORE] NARRATE: three paths watch the rack end. Top: the navigation camera feeds the learned proposer, which classifies the worker and predicts a heading; it sits above the proposal boundary and may only propose. Middle: the spinning lidar (the machine's is an Ouster OS1-64) feeds the clearance check of the permission path, which runs on its own safety microcontroller, as Lecture 4 placed it. That check is unrated. Bottom: a certified safety laser scanner, a planar time-of-flight device certified under IEC 61496-3, triggers the emergency stop when anything opaque enters its field; ISO 3691-4 requires person detection at Performance Level d, the second highest of five grades. It classifies nothing. ENGAGE: the camera also estimates depth; why must the clearance check use the lidar? (Expected: a checker that consumes the proposer's own perception shares every blind spot of the proposer; perception-side independence.) WARN: the scanner is the certified floor; Lecture 12 owns the certified floor beneath an unrated layer. FLEX: [CORE].}
\end{frame}

\begin{frame}{Vote: which range enters the clearance check?}
  \mlsysactivity{Vote}{2 min}
  \centering
  {\normalsize The rack end, a dock door open to daylight, a second robot in the aisle.}\\[0.35cm]
  \begin{minipage}{0.86\textwidth}\normalsize
    \choice{A}{a depth network on the navigation camera}\\[0.15cm]
    \choice{B}{an active stereo camera}\\[0.15cm]
    \answerat{2}{C}{the spinning lidar, qualified on this floor}\\[0.15cm]
    \choice{D}{an event camera}
  \end{minipage}\\[0.35cm]
  \uncover<2->{\small Each other choice has a failure this aisle supplies.}
  \mlsyscredit{Sensor Perception, \S8.2.3 to \S8.2.5 and table 1.}
  \note{[2 min] [CORE] ENGAGE: vote, discuss with a neighbor for 30 seconds, vote again. Step 2: C. A: a depth network infers scale from appearance; its spread is not measured error, and a photo printed on a box fools it; it also shares the proposer's camera. B: dock sunlight washes out the projected speckle and a neighbor's projector corrupts it; its error grows with range squared (segment 4). D: an event camera reports change, so a worker who stops at the corner disappears, and it gives no range at all. The lidar measures range directly and its error stays roughly flat with range. WARN: ``qualified on this floor'' matters, because specular floors and retroreflective vests still corrupt lidar (last segment). Land the segment aloud: sensing is chosen by the quantity the decision needs, and the range that may veto a proposal comes from a sensor the proposer does not use. FLEX: [CORE].}
\end{frame}

% =============================================================================
% SECTION: CAPTURE
% =============================================================================
\mlsysnote{[0.5 min] [CORE] A measurement of the right quantity is still useless until we know when it was taken and where the sensor was.}
\mlsysdivider{When and where was the measurement taken?}{Sensor Perception, \S8.2 and \S8.3}
\section{Capture}

\begin{frame}{A frame\textquoteright s capture epoch is the midpoint of its exposure}
  \centering
  \mlsyssteps[width=14.6cm,height=4.3cm,keepaspectratio]{capture_step}{3}\\[0.1cm]
  {\small Capture epoch: the instant the measurement represents.}\\[0.03cm]\uncover<2->{\small The arrival stamp pairs the frame with a pose \val{d_arrival} mm ahead of where the camera stood.}
  \mlsyscredit{Sensor Perception, \S8.2 and \S8.3; Napkin Math 8.2 stage values (illustrative).}
  \note{[2.5 min] [CORE] Step 1: the navigation camera integrates light for \val{t_exp} ms, a setting chosen for dim aisles. The image is dated at the exposure midpoint, the capture time $t_0$; the base moves \val{d_exposure} mm during the exposure itself. Step 2: the rows are then read out for \val{t_readout} ms, cross the link, and DMA completes \val{t_arrival} ms after $t_0$. A driver that stamps the frame when it arrives pairs it with a pose \val{d_arrival} mm farther down the aisle (\val{d_per_ms} mm per ms). Step 3: a hardware trigger marks the start of exposure, \val{t_trigger} ms before the midpoint; a driver that forgets to add half the exposure dates the frame \val{d_trigger} mm early. ENGAGE: auto-exposure changes the exposure every frame; what must each frame carry? (Expected: its own exposure time in metadata.) WARN: software stamps also inherit scheduler jitter; at a \val{yaw_rate} rad/s turn each millisecond misorients the camera by \val{misorient} mrad. The fix is a hardware stamp at capture. FLEX: [CORE].}
\end{frame}

\begin{frame}{A real log: one trigger, camera stamps \val{kitti_gap_whole} ms apart}
  \begin{columns}[c]
    \begin{column}{0.44\textwidth}
      {\small KITTI, sweep 20, host system clock}\par\vspace{0.1cm}
      \begin{logblock}[8.5]
        velodyne\_points/timestamps\_start.txt\\
        \hspace*{1em}13:02:27.755192959\\
        velodyne\_points/timestamps.txt\\
        \hspace*{1em}13:02:27.806744250\\
        image\_02/timestamps.txt (color)\\
        \hspace*{1em}13:02:27.817283584\\
        image\_00/timestamps.txt (gray)\\
        \hspace*{1em}13:02:27.823418112
      \end{logblock}
    \end{column}
    \begin{column}{0.54\textwidth}
      \centering\safeimg[width=\textwidth,height=4.3cm,keepaspectratio]{kitti_strip.pdf}
    \end{column}
  \end{columns}
  \vspace{0.1cm}
  {\small Across \val{kitti_sweeps} sweeps: color stamps +\val{kitti_color_off} ms, gray +\val{kitti_gray_off} ms after the trigger.}
  \mlsyscredit{KITTI raw, drive 2011\_09\_26\_0001 extract, timestamp files (excerpt, date omitted); Geiger et al., IJRR 2013, \S III.A. Medians: calcs.py.}
  \note{[2 min] [OPTIONAL] Real timestamp files from KITTI. The paper says a reed contact on the spinning lidar triggers all four cameras when it faces forward, and that every timestamp was recorded on the host computer's system clock. NARRATE: the lidar's forward stamp is \val{kitti_fwd} ms into this \val{kitti_sweep_ms} ms sweep. The two color cameras' stamps land \val{kitti_color_off} ms after the trigger and the two gray cameras' \val{kitti_gray_off} ms after (medians over \val{kitti_sweeps} sweeps). The light reached all four sensors on one trigger, so the \val{kitti_gap} ms between them is host-side delay, our reading of the log. ENGAGE: a dataset user pairs gray images with lidar by the stamps. How much does the car move in \val{kitti_gray_off} ms at its \val{kitti_v} m/s? (Expected: about \val{kitti_gray_travel_cm} cm; estimate aloud.) WARN: KITTI documents this and its sync data handle it; the lesson is that an arrival stamp silently carries delay. FLEX: [OPTIONAL].}
\end{frame}

\begin{frame}{Rows read in sequence shear a moving scene}
  \mlsysactivity{Estimate}{2.5 min}
  \begin{columns}[c]
    \begin{column}{0.27\textwidth}
      \centering\safeimg[width=\textwidth,height=3.0cm,keepaspectratio]{propeller-rolling-shutter.jpg}\\[0.05cm]
      {\footnotesize A propeller, one phone frame}
    \end{column}
    \begin{column}{0.38\textwidth}
      \begin{mlsyscard}{mlsysaccent}
        \small\textbf{Given}\\[0.05cm]
\val{rows} rows in \val{t_readout} ms\\
        \val{v_drive} m/s along the aisle\\
        Yaw \val{yaw_rate} rad/s at \val{d_clear} m range\\
        Allowance \val{delta_loc} mm
      \end{mlsyscard}
      \small
      \uncover<2->{Row interval: \textbf{\val{t_row} \textmu s}\\
      Along: $\val{v_drive} \times \val{t_readout} = \mathbf{\val{d_skew}}$ mm\\
      Sideways: \textbf{\val{d_lateral}} mm}
    \end{column}
    \begin{column}{0.33\textwidth}
      \centering\uncover<2->{\safeimg[width=\textwidth,height=4.6cm,keepaspectratio]{rows.pdf}}
    \end{column}
  \end{columns}
  \mlsyscredit{Photo: Dick Lyon, CC BY-SA 4.0, Wikimedia Commons. Sensor Perception, Napkin Math 8.1 and equation 1 (rows: datasheet; speed and yaw: illustrative).}
  \note{[2.5 min] [CORE] The photo: a rolling-shutter phone camera reads rows top to bottom while the blades turn, so they come out bent and detached. The navigation camera does the same. ENGAGE: two minutes in pairs with the card: how far apart, along the aisle, do the top and bottom rows see a rack upright? Step 2: each row is read \val{t_readout} ms / \val{rows} = \val{t_row} \textmu s after the one above. The last row sees the upright from \val{v_drive} m/s $\times$ \val{t_readout} ms = \val{d_skew} mm farther down the aisle, half of the \val{delta_loc} mm localization allowance. Turning into the aisle adds \val{d_lateral} mm sideways at \val{d_clear} m, \val{lateral_pct} percent of the \val{side_clearance} mm side clearance. At \val{v_aisle} m/s the skew is still \val{d_skew_aisle} mm. WARN: slowing down barely helps; the fix is to unwarp each row with the pose interpolated from the \val{f_imu} Hz IMU before stereo or detection runs. A global-shutter camera, like the wrist camera, removes the skew at a higher cost. FLEX: [CORE].}
\end{frame}

\begin{frame}{A lidar sweep takes \val{kitti_sweep0} ms; the car moved \val{kitti_sweep0_travel} m}
  \begin{columns}[c]
    \begin{column}{0.52\textwidth}
      \centering\safeimg[width=\textwidth,height=5.0cm,keepaspectratio]{kitti_sweep.pdf}
    \end{column}
    \begin{column}{0.46\textwidth}
      \centering\safeimg[width=\textwidth,height=3.0cm,keepaspectratio]{ouster-os1-cloud.jpg}\\[0.05cm]
      {\footnotesize Ouster OS1-64 returns over 18 s, each placed with the pose at its own firing time.}\\[0.2cm]
      {\small Our machine at \val{f_lidar} Hz: \val{t_sweep} ms, \textbf{\val{d_sweep} mm} per sweep.}
    \end{column}
  \end{columns}
  \mlsyscredit{Left: KITTI sync sweep 0 and OXTS speed (CC BY-NC-SA 3.0); time from trigger by angle, ours. Right: Daniel L. Lu, CC BY 4.0, Wikimedia Commons. Sensor Perception, \S8.2.5.}
  \note{[2 min] [OPTIONAL] Left: one real KITTI sweep from above, the car at the center driving at \val{kitti_v} m/s (\val{kitti_kmh} km/h, from its GPS/IMU log). The sweep starts and ends facing backward and passes the front, where it triggers the cameras, halfway. Color shows each return's time from the trigger, assigned by angle. Returns just left and right of the seam behind the car were fired \val{kitti_sweep0} ms apart, and the car moved \val{kitti_sweep0_travel} m between them. NARRATE: a spinning lidar is a rolling shutter turned on its side. Our machine's OS1-64 at \val{f_lidar} Hz takes \val{t_sweep} ms per sweep, \val{d_sweep} mm of travel and \val{rot_sweep} mrad of turn at the drive limit. De-skewing moves each return to the frame's epoch using the pose at its firing time; the right image keeps crisp facades only because of that. ENGAGE: why stamp every column instead of the whole sweep? (Expected: a single sweep time is wrong by up to \val{t_sweep} ms for some points.) FLEX: [OPTIONAL].}
\end{frame}

\begin{frame}{Every claim reaches the world through a chain of frames}
  \begin{columns}[c]
    \begin{column}{0.40\textwidth}
      \centering\safeimg[width=\textwidth,height=5.5cm,keepaspectratio]{frames.pdf}
    \end{column}
    \begin{column}{0.58\textwidth}
      {\small A real calibration record (KITTI, 2011\_09\_26)}\par\vspace{0.08cm}
      \begin{logblock}[8]
        calib\_velo\_to\_cam.txt\\
        calib\_time: 15-Mar-2012 11:37:16\\
        R: 7.533745e-03 -9.999714e-01 ...\\
        T: -4.069766e-03 -7.631618e-02 -2.717806e-01\\[0.15cm]
        calib\_cam\_to\_cam.txt\\
        K\_02: 9.597910e+02 0.000000e+00 6.960217e+02 ...
      \end{logblock}
      \vspace{0.2cm}
      {\small A \val{angle_err}° yaw error moves a point \textbf{\val{d_deg_rack} mm} at \val{d_clear} m and \textbf{\val{d_deg_far} mm} at \val{mount_range} m.}
    \end{column}
  \end{columns}
  \mlsyscredit{Frames after REP 105 and Sensor Perception figure 4 and table 2. KITTI calibration files (excerpt; values truncated at ``...''). Mount error: Sensor Perception, \S8.2.2 and \S8.3.}
  \note{[2.5 min] [CORE] Left, after the chapter's frame tree and the REP 105 convention of the Robot Operating System: map may jump when the robot relocalizes, odom stays continuous but drifts, base\_link is the body, and each sensor\_link hangs from it by a calibrated mount. An observation is placed in the world by composing every transform at its capture time. Right: a real calibration record, the rotation R and translation T from KITTI's lidar to its camera, and the color camera's intrinsic matrix K (focal lengths and principal point). NARRATE: angular errors grow with range. A \val{angle_err}° mount error moves a point \val{d_deg_rack} mm at the rack end and \val{d_deg_far} mm at \val{mount_range} m, \val{mount_over_side} times the side clearance. ENGAGE: mounts loosen and links warm; a \val{link_len} m aluminum link grows \val{d_link} mm for \val{link_dt} K. What must every observation carry? (Expected: the calibration revision it used, with its validity range.) LINK: Lecture 9 owns how these transforms go stale as joints move. FLEX: [CORE].}
\end{frame}

\begin{frame}{Timing slips cost tens of millimeters; a loose mount costs more}
  \centering
  \safeimg[width=14.6cm,height=4.6cm,keepaspectratio]{errbars.pdf}\\[0.1cm]
  \small Each error is invisible in the image and visible only in the record.
  \mlsyscredit{Sensor Perception, \S8.3 figure (values from calcs.py at the \val{v_drive} m/s drive limit).}
  \note{[1.5 min] [CORE] Land the segment with the chapter's figure. Every bar is a position error the machine would carry silently. The trigger stamp without half the exposure: \val{d_trigger} mm. Row skew: \val{d_skew} mm. The arrival stamp: \val{d_arrival} mm. The dispatch age, next segment: \val{d_age} mm. A \val{angle_err}° mount error read at \val{mount_range} m down a long aisle: \val{d_deg_far} mm, past the \val{side_clearance} mm side clearance, so the planner either steers into the racking or refuses an open aisle. ENGAGE: which of these does a better neural network fix? (Expected: none; each is fixed by a hardware stamp, row compensation, or calibration.) Land aloud: the machine can use a measurement only when it is dated at capture and placed through calibrated frames. FLEX: [CORE].}
\end{frame}

% =============================================================================
% SECTION: AGE
% =============================================================================
\mlsysnote{[0.5 min] [CORE] Dated and placed, the observation is still old by the time anyone reads it. How old, and what does that cost?}
\mlsysdivider{How old is the observation when it is used?}{Sensor Perception, \S8.4 and \S8.5}
\section{Age}

\begin{frame}{Predict: where do the \val{t_age} ms go?}
  \mlsysactivity{Predict}{3 min}
  \centering
  \mlsyssteps[width=14.6cm,height=5.0cm,keepaspectratio]{waterfall_step}{4}\\[0.05cm]
  \uncover<4->{\small Observation age at dispatch: \textbf{\val{t_age} ms}, \textbf{\val{d_age} mm} of travel before any consumer reads the frame.}
  \mlsyscredit{Sensor Perception, Napkin Math 8.2 and table 3: stage values are budgeted allocations (readout: datasheet; others illustrative).}
  \note{[3 min] [CORE] Observation age runs from the exposure midpoint to the moment perception dispatches the observation. Step 1: the total is \val{t_age} ms. ENGAGE: vote: which is larger, the vision backbone or everything before any computation starts? Step 2: half the exposure, \val{t_st_exphalf} ms, plus the rolling readout, \val{t_st_readout} ms: \val{t_capture} ms before a single kernel runs. Step 3: link \val{t_st_link}, DMA and queue \val{t_st_dma}, image signal processor \val{t_st_isp}: \val{t_through_isp} ms. Step 4: the INT8 backbone, \val{t_st_backbone} ms, and dispatch \val{t_st_ipc} ms: \val{t_age} ms, \val{d_age} mm at \val{v_drive} m/s (calcs.py). Students who expected the model to dominate learn that it is \val{pct_backbone} percent; capture is \val{pct_capture} percent. WARN: each value is an allocation the stage must meet at its tail, and the last column of table 3 names the measurement that must confirm it. LINK: Lecture 2 set the freshness deadline a measurement must meet; Lecture 4 allotted the sensing handoff; this is the age itself. FLEX: [CORE].}
\end{frame}

\begin{frame}{A measured detector runs in \val{rn_p99_short} ms; capture took \val{t_capture} ms}
  \begin{columns}[c]
    \begin{column}{0.48\textwidth}
      \begin{logblock}[8]
        MLPerf Results Summary\\
        SUT name : LWIS\_Server\\
        Scenario : SingleStream\\
        Min latency (ns)\hfill 11410286\\
        50.00 percentile latency (ns)\hfill 11560481\\
        \textcolor{mlsysaccent}{99.00 percentile latency (ns)\hfill 11790504}\\
        Max latency (ns)\hfill 11836933
      \end{logblock}
      \vspace{0.1cm}
      \centering\safeimg[width=0.62\textwidth,height=1.9cm,keepaspectratio]{jetson-agx-orin.png}\\
      {\footnotesize NVIDIA Jetson AGX Orin}
    \end{column}
    \begin{column}{0.50\textwidth}
      \centering\safeimg[width=\textwidth,height=4.6cm,keepaspectratio]{mlperf.pdf}
    \end{column}
  \end{columns}
  \mlsyscredit{MLPerf Inference v3.1, closed, NVIDIA Orin\_TRT (Jetson AGX Orin 64G), RetinaNet, INT8, SingleStream, run 1, mlperf\_log\_summary.txt (excerpt). Photo: Auledas, CC BY 4.0, Wikimedia Commons.}
  \note{[1.5 min] [OPTIONAL] Real tool output: MLPerf's log for the RetinaNet object detector, INT8 weights, on a Jetson AGX Orin, one image at a time. The 99th percentile is \val{rn_p99} ms. NARRATE: the chapter's backbone allocation is \val{t_st_backbone} ms, \val{backbone_over_rn} times this measured detector; capture before computation is \val{t_capture} ms, \val{capture_over_rn} times it. A faster model shortens one bar of seven. WARN: MLPerf starts its clock with the image already in memory, so this number is the backbone stage alone; it excludes exposure, readout, link, DMA, and the image processor. A different model on a different input would differ; the point is the scale. ENGAGE: what would halve the \val{t_capture} ms? (Expected: a shorter exposure in brighter light, a faster readout mode, or a global shutter with a known row convention.) FLEX: [OPTIONAL].}
\end{frame}

\begin{frame}{Observation age is the fifth row of the stopping budget}
  \centering
  \only<1|handout:0>{\mlsyskitfig[width=0.95\textwidth,height=4.4cm,keepaspectratio]{budget_bar_step4}}%
  \only<2-|handout:1>{\mlsyskitfig[width=0.95\textwidth,height=4.4cm,keepaspectratio]{budget_bar_step5}}\par\vspace{0.05cm}
  \uncover<3->{\small The \val{d_spare15} mm left must cover admission delay, a jerk-limited stop, and a tracking inset.}
  \mlsyscredit{Sensor Perception, \S8.4.1; kit bar; sums: calcs.py.}
  \note{[2.5 min] [CORE] Step 1: the bar from the recall; braking, allowances, onset, and lease path come to \val{d_rows14} mm. Step 2: this lecture's segment, the observation age, speed times age: \val{v_drive} m/s $\times$ \val{t_age} ms = \val{d_age} mm. It was always inside Lecture 1's \val{t_tau} ms delay, so the bar does not lengthen; it names the piece: \val{t_age} + \val{t_rest} = \val{t_tau} ms, and the partial stopping distance stays \val{d_rows15} mm. Step 3: \val{d_spare15} mm of the \val{d_clear} m remain, and that is not spare. It must absorb the travel between dispatch and admission, which Lecture 13 allocates; the longer stop of a jerk-limited brake, which Lecture 11 owns; and the tracking inset of the stopping envelope, which Lecture 12 owns. ENGAGE: the permission path never reads the image; can it shorten this row? (Expected: no; the age reaches it through the proposal's stamped capture time, and pose propagation can place the old observation correctly but cannot return distance already traveled.) FLEX: [CORE].}
\end{frame}

\begin{frame}{Evidence keeps aging after dispatch}
  \centering
  \safeimg[width=14.6cm,height=3.9cm,keepaspectratio]{agegrow.pdf}\\[0.15cm]
  \small Each consumer computes age on its own clock from the capture time.\\[0.05cm]
  \uncover<2->{\small One dropped frame: \val{t_frame} ms older, \textbf{\val{d_frame} mm}, \val{frame_pct}\% of the \val{d_spare15} mm.}
  \mlsyscredit{Sensor Perception, Definition 8.1, \S8.8 (dispatch threshold \val{t_age} + \val{t_sync} ms) and Checkpoint 8.1. Values: calcs.py.}
  \note{[2 min] [CORE] Observation age stops at dispatch; the evidence does not. The current evidence age is now minus the capture time, read on the consumer's own monotonic clock, and it keeps growing while the observation waits for admission and while it is reused tick after tick. The record therefore stores the capture time, never an age. The producer refuses to dispatch a frame older than \val{t_dispatch_limit} ms: the \val{t_age} ms budget plus the \val{t_sync} ms clock-conversion bound. The permission path applies a separate, later limit at admission. ENGAGE: step 2: if a dropped frame leaves a proposal resting on the previous frame, one period older at \val{f_frame} Hz, what does it cost? (Expected: \val{t_frame} ms, \val{d_frame} mm, \val{frame_pct} percent of the remainder.) WARN: a new proposal built from a buffered frame does not refresh the evidence; its age runs from the old capture. FLEX: [CORE].}
\end{frame}

\begin{frame}{Diagnose: why do the gray-camera stamps wander?}
  \mlsysactivity{Diagnose}{2.5 min}
  \begin{columns}[c]
    \begin{column}{0.50\textwidth}
      \centering\safeimg[width=\textwidth,height=4.7cm,keepaspectratio]{kitti_jitter.pdf}
    \end{column}
    \begin{column}{0.48\textwidth}
      \small One trigger fires every camera every \val{kitti_sweep_ms} ms.\\[0.15cm]
      \choice{A}{the gray cameras expose longer}\\[0.08cm]
      \answerat{2}{B}{host delays enter the stamps}\\[0.08cm]
      \choice{C}{the lidar spins unevenly}\\[0.2cm]
      \uncover<2->{\footnotesize On our machine the same queue tail ages frames past \val{t_dispatch_limit} ms; batching \val{batch_words} frames adds \val{t_batch_wait} ms.}
    \end{column}
  \end{columns}
  \mlsyscredit{KITTI raw data, drive 2011\_09\_26\_0001 extract, all timestamp files (CC BY-NC-SA 3.0); periods: calcs.py. Batching: Sensor Perception, \S8.4.3 and \S8.5.}
  \note{[2.5 min] [CORE] Real data: the gap between consecutive stamps in the KITTI log, every frame plotted. The lidar holds \val{kitti_per_lidar_min} to \val{kitti_per_lidar_max} ms; the color cameras \val{kitti_per_color_min} to \val{kitti_per_color_max}; the gray cameras wander from \val{kitti_per_gray_min} to \val{kitti_per_gray_max}. ENGAGE: one minute; the trigger is the lidar. Step 2: B. The light arrives on schedule; the stamps record when the host handled each frame, so queueing and scheduling on the host leak into the timeline. A is wrong because exposure would shift every stamp equally; C is ruled out by the lidar's own steady period. NARRATE the machine's version: camera DMA bursts and inference share one memory system on an application processor like the Orin, so a frame can wait behind traffic and reach dispatch older than its budget; the contract refuses it. Batching \val{batch_words} frames at \val{f_frame} Hz holds the first one \val{t_batch_wait} ms, \val{t_batched} ms in all. LINK: Lecture 13 owns the contention mechanism and its placement remedies. FLEX: [CORE].}
\end{frame}

\begin{frame}{Delay spends phase margin; at \val{t_pm} ms none remains}
  \begin{columns}[c]
    \begin{column}{0.56\textwidth}
      \centering\safeimg[width=\textwidth,height=4.6cm,keepaspectratio]{phase.pdf}
    \end{column}
    \begin{column}{0.42\textwidth}
      \small A visual-servo loop crosses over at \val{wc} rad/s with \val{pm_deg}° of margin.\\[0.2cm]
      $t_{\text{max}} = \dfrac{\val{pm_rad}\ \text{rad}}{\val{wc}\ \text{rad/s}} = \val{t_pm}$ ms\\[0.2cm]
      \val{t_pm_used} ms of perception spends \textbf{\val{pm_used_pct}\%} of it.
    \end{column}
  \end{columns}
  \mlsyscredit{Sensor Perception, \S8.4.4 and figure (illustrative loop).}
  \note{[1 min] [OPTIONAL] The same age has a second cost. A pure delay removes phase in proportion to frequency, so a loop closed through perception loses margin as age grows. With \val{pm_deg} degrees (\val{pm_rad} rad) at a \val{wc} rad/s crossover, about \val{f_c} Hz, the margin is gone at \val{t_pm} ms; \val{t_pm_used} ms of exposure, transport, and inference spends \val{pm_used_pct} percent of it. LINK: Lecture 4 owns ``delay spends phase margin''; this is its perception term. WARN: the two costs are paid in different currencies, travel against clearance and phase against stability, so the age is checked against each and counted once in each. FLEX: [OPTIONAL]. Land the segment aloud: observation age is paid in travel before any consumer can react, and no later stage returns it.}
\end{frame}

% =============================================================================
% SECTION: ERROR
% =============================================================================
\mlsysnote{[0.5 min] [CORE] A fresh, placed measurement still has error. How large, measured how, and charged to what?}
\mlsysdivider{How far can the measurement be trusted?}{Sensor Perception, \S8.6}
\section{Error}

\begin{frame}{Estimate: where does stereo meet its range wall?}
  \mlsysactivity{Estimate}{3 min}
  \begin{columns}[c]
    \begin{column}{0.36\textwidth}
      \begin{mlsyscard}{mlsysaccent}
        \small\textbf{Given: a stereo pair}\\[0.05cm]
        Baseline \val{st_b} m, focal \val{st_f} px\\
        Disparity error \val{st_sd} px\\
        $\sigma_z = \dfrac{z^2\,\sigma_d}{f\,b}$
      \end{mlsyscard}
      \small
      \wstep{3}{$\sigma_z = \val{delta_loc}$ mm at \textbf{\val{st_z_loc} m}}
      \wstep{3}{$\sigma_z = \val{side_clearance}$ mm at \textbf{\val{st_z_side} m}}
    \end{column}
    \begin{column}{0.62\textwidth}
      \centering\mlsyssteps[width=\textwidth,height=5.3cm,keepaspectratio]{stereo_step}{3}
    \end{column}
  \end{columns}
  \mlsyscredit{Sensor Perception, Napkin Math 8.3 and equation 8 (illustrative rig; the machine carries no depth camera). Ouster OS1 datasheet, rev7: precision \val{os1_prec_min_cm} to \val{os1_prec_max_cm} cm, 1$\sigma$, beyond 1 m.}
  \note{[3 min] [CORE] Stereo recovers depth from disparity, $z = f b / d$; differentiating gives the error card's formula, so the error grows with range squared. Step 1: at the rack end, \val{d_clear} m, one standard deviation is \val{st_rack} mm; fine. ENGAGE: two minutes: at what range does one sigma alone use the \val{delta_loc} mm localization allowance, and the \val{side_clearance} mm side clearance? Step 2: the lidar's datasheet precision, flat between \val{os1_prec_min} and \val{os1_prec_max} mm, because time of flight does not triangulate. Step 3: $z = \sqrt{C f b / \sigma_d}$: \val{st_z_loc} m and \val{st_z_side} m. \Val{st_range_ratio_words} times the range costs \val{st_growth} times the error (\val{st_sigma_near} mm at \val{st_near} m, \val{st_sigma_far} at \val{st_far} m). Halving the baseline pulls the \val{st_z_side} m wall in to \val{st_z_side_half} m. WARN: these are one-sigma scales and bound nothing; mismatches on corrugated siding or shiny floors give gross outliers far outside the Gaussian. FLEX: [CORE].}
\end{frame}

\begin{frame}{Two numbers in the contract must be measured on this floor}
  \begin{columns}[c]
    \begin{column}{0.56\textwidth}
      {\small Ouster OS1 datasheet, rev7 (excerpt)}\par\vspace{0.08cm}
      \begin{logblock}[8.5]
        Range Precision\hfill Min: ±0.5 cm, Max: ±3 cm\\
        Range Accuracy\hfill ±2.5 cm for lambertian targets,\\
        \hfill ±5 cm for retroreflective targets\\
        Data Latency\hfill < 10 ms\\
        IEEE1588 Precision Time Protocol (PTP);\\
        \hfill Accuracy: <1 ms error\\
        Internal 10 ppm drift clock
      \end{logblock}
    \end{column}
    \begin{column}{0.42\textwidth}
      \small\textbf{\textcolor{figveto}{1. The residual tail}} on this floor, this mode, this motion.\\[0.3cm]
      \textbf{\textcolor{figamber}{2. The clock-conversion bound}} from the lidar's clock to the safety microcontroller's: offset tail, path asymmetry, drift.
    \end{column}
  \end{columns}
  \mlsyscredit{Ouster, Inc., OS1 datasheet, hardware rev7, firmware 3.1 (excerpt). Sensor Perception, \S8.6.2.}
  \note{[2 min] [CORE] Real datasheet lines. Precision is the spread of 100 measurements on a static target; accuracy is the bias of their mean, \val{os1_acc} mm on diffuse surfaces, a calibration offset that averaging cannot remove, and twice that on retroreflectors such as safety vests. NARRATE: neither number was taken on our floor, in our scan mode, while the base moves. Two values in the claim must come from a measurement campaign. First, the tail of the residuals against ground truth, which sets the clearance inset (next slide). Second, the clock-conversion bound: every free-running oscillator is its own clock domain, and converting a capture time into the safety microcontroller's clock carries error. PTP assumes equal path delays both ways; an asymmetry shifts the offset by half of itself with no symptom. A \val{os1_drift_ppm} ppm clock drifts \val{drift_us} \textmu s per second between corrections. The datasheet promises only under \val{os1_ptp} ms. FLEX: [CORE].}
\end{frame}

\begin{frame}{The clearance inset rests on a measured 99th-percentile residual}
  \begin{columns}[c]
    \begin{column}{0.57\textwidth}
      \centering\mlsyssteps[width=\textwidth,height=4.9cm,keepaspectratio]{inset_step}{2}
    \end{column}
    \begin{column}{0.41\textwidth}
      \small $C_{\text{inset}} = C_{\text{safe}} + \lVert e\rVert_{\text{p99}}$\\[0.25cm]
      \uncover<2->{Sort $n$ residuals; take rank $\lceil (n+1)(1-\alpha) \rceil$.\\[0.1cm]
      $\alpha = \val{alpha_pct}\%$: at least \textbf{\val{n_min}} residuals;\\ of \val{n_resid}, the \textbf{\val{rank}st}.\\[0.15cm]
      {\footnotesize A new floor or a recalibration needs a new campaign.}}
    \end{column}
  \end{columns}
  \mlsyscredit{Sensor Perception, equation 9, Definition 8.2, \S8.6.2 (split-conformal rank, Angelopoulos and Bates 2023). Distribution schematic.}
  \note{[2.5 min] [CORE] Step 1: the permission path holds the estimated obstacle at $C_{\text{inset}}$, so the true obstacle stays beyond the required clearance $C_{\text{safe}}$ except in the tail. Step 2: the inset is the required clearance plus a 99th-percentile quantile of the full residual norm, measured against ground truth, so it contains the bias and outliers a covariance omits. The split-conformal rule picks the sorted residual of rank $\lceil (n+1)(1-\alpha) \rceil$: for \val{alpha_pct} percent it needs at least \val{n_min} residuals and returns the largest; with \val{n_resid} it returns the \val{rank}st, the \val{rank_top}th largest. It holds if future residuals are exchangeable with the campaign's. ENGAGE: the shortcut uses one stereo sigma: \val{inset_near} mm at \val{st_near} m, \val{inset_far} mm at \val{st_far} m. What does it miss? (Expected: everything beyond one sigma and every mismatch.) LINK: Lecture 10's conformal threshold is the same construction, applied there to model scores. WARN: a 99th percentile leaves one residual in a hundred above it; that tail needs a stop the permission path can still execute. FLEX: [CORE].}
\end{frame}

\begin{frame}{Qualified lidar range holds a metric veto over a learned label}
  \centering
  \safeimg[width=14.6cm,height=4.4cm,keepaspectratio]{veto.pdf}\\[0.15cm]
  \small A \val{f_vlm} Hz model says ``clear aisle''; \val{sweeps_per_vlm_words} lidar sweeps arrive before it speaks again.
  \mlsyscredit{Sensor Perception, Systems Perspective 8.3 and \S8.6.3 (illustrative timeline; \val{f_vlm} Hz and \val{f_lidar} Hz rates from the chapter).}
  \note{[2 min] [CORE] The metric veto. Top: a slow vision-language model labels the aisle ``clear'' twice a second with high confidence. Middle: the lidar sweeps \val{f_lidar} times a second. At about 650 ms a return appears inside the stopping envelope. Bottom: the permission path halts at its next check, whatever the label says. NARRATE: learned depth and semantic confidence are uncalibrated scores; a depth network can err by meters on a new floor texture or a life-size photo printed on a box, and its predicted spread was never measured against ground truth. So a learned head may classify and propose, while the range that enters the clearance check comes from qualified metric sensing. ENGAGE: why would the model's label still be useful? (Expected: it tells a person from a tote, sets intent and heading; it never erases an occupied return.) LINK: Lecture 3 owns the softmax illusion. Land the segment: claimed error must be measured, and metric evidence outranks semantic confidence. FLEX: [CORE].}
\end{frame}

% =============================================================================
% SECTION: CONTRACT
% =============================================================================
\mlsysnote{[0.5 min] [CORE] Last question: what fails without raising a flag, and what must every observation carry so a consumer can refuse it?}
\mlsysdivider{What must a record say before anyone reads it?}{Sensor Perception, \S8.7 and \S8.8}
\section{Contract}

\begin{frame}{Glare clips pixels in frames that arrive on time}
  \centering
  \safeimg[width=14.6cm,height=4.5cm,keepaspectratio]{kitti_sat.pdf}\\[0.15cm]
  \small A histogram monitor sees the clipping; the driver reports a healthy frame.
  \mlsyscredit{KITTI raw data, drive 2011\_09\_26\_0001 sync, image\_02 frame 0 (CC BY-NC-SA 3.0); clipped share: calcs.py. Degradation modes: Sensor Perception, \S8.7.1.}
  \note{[2 min] [OPTIONAL] Real frame: \val{sat_pct} percent of the pixels are clipped at \val{pix_max}, here the sky. A pixel-intensity histogram right after DMA flags that; the driver, the link, and the checksum all report a perfect frame at the nominal rate. Now turn the robot toward a dock door, and the clipped patch hides a rack. NARRATE the other silent modes by their physics: fog makes contrast decay exponentially (at \val{fog_beta} per meter it falls below \val{fog_eps} by \val{d_fog} m); fog droplets backscatter lidar into a false wall 0.5 to 3 m out; retroreflective vests make returns too strong, so they read short; polished or wet floors bounce the beam away (``puddle black holes'') or create phantom points behind glass. Input-support monitors set bits such as BIT\_OPTICAL\_SATURATED or BIT\_OOD\_DETECTED in a health mask, and the camera can be demoted to lidar-only operation. LINK: Lecture 7 owns the runtime monitor channels. FLEX: [OPTIONAL].}
\end{frame}

\begin{frame}{Vote: a lost frame or a misdated one?}
  \mlsysactivity{Vote}{2.5 min}
  \centering
  {\normalsize Which hurts the estimate more: dropping a frame, or accepting it \val{t_arrival} ms misdated?}\\[0.15cm]
  \uncover<2->{\safeimg[width=14.6cm,height=4.3cm,keepaspectratio]{lostmis.pdf}}\\[0.05cm]
  \uncover<2->{\small $\Delta x = v\,\delta t + \tfrac{1}{2} a\,\delta t^2$: a wrong stamp displaces the obstacle and narrows the spread.}
  \mlsyscredit{Sensor Perception, \S8.7 and equation 11 (schematic distributions; offset from calcs.py).}
  \note{[2.5 min] [CORE] ENGAGE: vote, then a neighbor, then vote again. Step 2: the misdated frame. Left: a lost frame raises a flag; the estimator gets no update, its spread grows under its motion model, and the evidence is honestly one frame, \val{t_frame} ms, older. Right: a frame stamped at arrival is accepted as information, so the estimator narrows its spread around a position \val{d_arrival} mm from where the obstacle stood (equation 11 at constant speed). The machine grows more confident as it grows more wrong. WARN: students rank faults by how loud they are; here the quiet one is worse. Detectable faults (lost link, CRC failure, buffer overrun) age out honestly; silent ones (saturation, smudged lens, fog backscatter, a misdated frame) deliver valid tensors. The design goal is to turn silent modes into detectable ones. LINK: Lecture 9 owns the Kalman filter in which this happens. FLEX: [CORE].}
\end{frame}

\begin{frame}{True or false? Two claims about sensor evidence}
  \mlsysactivity{Vote}{1.5 min}
  \begin{enumerate}\setlength{\itemsep}{0.3cm}
    \item \normalsize A detection with 99\% confidence may be admitted whatever its age.\\
      \uncover<2->{\small\textbf{\textcolor{figveto}{False.}} \Val{held_frames_words} frames late is \val{t_batch_wait} ms, \val{d_batch_wait} mm, and a worker who stepped out after exposure is not in the image.}
    \item \normalsize Lidar projects its own light, so mirrors and puddles cannot fool it.\\
      \uncover<3->{\small\textbf{\textcolor{figveto}{False.}} Specular floors drop returns; glass makes phantoms; vests read short (±\val{os1_acc_retro_cm} cm datasheet accuracy).}
  \end{enumerate}
  \mlsyscredit{Sensor Perception, \S8.9 (fallacies 1 and 6); Ouster OS1 datasheet, rev7. Travel: calcs.py.}
  \note{[1.5 min] [OPTIONAL] Two of the chapter's fallacies as a quick vote, hands up for true. Step 2: false. Confidence describes the classifier's view of a past image. \Val{held_frames_words} frames at \val{f_frame} Hz is \val{t_batch_wait} ms, \val{d_batch_wait} mm of travel at the drive limit; the observation must pass its freshness check, and state propagated without fresh data must widen its uncertainty. Step 3: false. Active sensors assume diffuse surfaces. Polished or wet concrete reflects the beam away and leaves holes in the cloud; glass and metal produce multi-bounce phantoms behind walls or under the floor; retroreflectors trigger the detector early and bloom into neighboring channels, which is why the OS1 datasheet doubles its accuracy tolerance for them. FLEX: [OPTIONAL]; cut if short and mention both in one sentence each.}
\end{frame}

\begin{frame}{The observation contract lets a consumer refuse before reading}
  \centering{\small Record 6 of the record chain. The payload is read last.}\\[0.15cm]
  \small\renewcommand{\arraystretch}{0.75}
  \begin{tabular}{@{}P{3.6cm}P{5.2cm}P{5.2cm}@{}}\toprule
    \textbf{Field} & \textbf{Navigation camera} & \textbf{Consumer check} \\\midrule
    \urow{1}{capture time, clock}{exposure midpoint, cam\_nav}{age on its own clock}
    \urow{1}{clock bound}{\val{t_sync} ms: dispatch limit \val{t_dispatch_limit} ms}{refuse if missing}
    \urow{2}{frame, calibration}{nav\_cam, revision}{reject outside valid range}
    \urow{3}{covariance, inset}{inset from p99 residual}{inflate the occupied region}
    \urow{4}{health, sequence}{saturation, sync, OOD bits}{refuse on critical bits}
    \urow{4}{CRC, build digest}{CRC-32, SHA-256}{integrity only}\bottomrule
  \end{tabular}
  \mlsyscredit{Sensor Perception, \S8.8, equation 12 and table 5 (navigation camera values; clock bound illustrative).}
  \note{[3 min] [CORE] The record that wraps each observation, built by row groups. Step 1, time: the exposure-midpoint capture time on a named monotonic clock, and the bound on converting it to the safety microcontroller's clock, \val{t_sync} ms here, so the producer refuses to dispatch past \val{t_age} + \val{t_sync} = \val{t_dispatch_limit} ms. No age is stored; each consumer computes it. Step 2, place: the frame id and the calibration revision with its validity range. Step 3, error: the covariance and the clearance inset from measured residuals. Step 4, health and integrity: the health bits, a sequence counter that exposes drops and reordering, a CRC against accidental corruption, and a digest of the firmware and model build. ENGAGE: in what order should a consumer check these? (Expected: integrity and health, clock conversion and freshness, frame validity, then error, and only then the payload.) WARN: a CRC is not authentication and a clean record is not semantic truth. LINK: Lecture 4's handoff record fixed the shared clock; a fixed schema also lets Lecture 15 inject faults field by field. FLEX: [CORE].}
\end{frame}

\begin{frame}{A ROS 2 header carries a time and a frame, without their limits}
  \begin{columns}[c]
    \begin{column}{0.54\textwidth}
      {\small std\_msgs/msg/Header.msg (ROS 2 Humble)}\par\vspace{0.05cm}
      \begin{logblock}[8.5]
        builtin\_interfaces/Time stamp\\[0.05cm]
        \# Transform frame with which this data is associated.\\
        string frame\_id
      \end{logblock}
      \vspace{0.15cm}
      {\small sensor\_msgs/msg/Imu.msg (excerpt)}\par\vspace{0.05cm}
      \begin{logblock}[8.5]
        \# A covariance matrix of all zeros will be interpreted as "covariance unknown", ...
      \end{logblock}
    \end{column}
    \begin{column}{0.44\textwidth}
      \small Missing from the header:\\[0.1cm]
      clock id and conversion bound\\
      calibration revision\\
      measured inset\\
      health bits
    \end{column}
  \end{columns}
  \mlsyscredit{ROS 2 common\_interfaces, humble branch: std\_msgs/msg/Header.msg and sensor\_msgs/msg/Imu.msg (exact lines; Apache-2.0).}
  \note{[1 min] [OPTIONAL] Real source: the header every ROS 2 sensor message carries. It has the two fields the chapter calls mandatory first, a stamp and a frame id, and nothing that says which clock produced the stamp, how well it converts, which calibration the frame used, or whether the sensor was healthy. The IMU message shows the usual convention for error: a covariance of all zeros means unknown. NARRATE: the observation contract does not replace ROS messages, Protocol Buffers, or FlatBuffers; any of them can carry its fields. It adds checkable physical invariants on top of a transport schema. ENGAGE: which missing field would have caught KITTI's 6 ms gap between cameras? (Expected: a clock id with a hardware capture stamp.) FLEX: [OPTIONAL].}
\end{frame}

\mlsysnote{[1 min] [CORE] The chapter's invariant, read aloud. Pass a claim when its capture time and clock domain, its coordinate frame, and its qualified uncertainty can all be checked for the proposed motion; refuse it when any cannot. This is the third law, stated today, at the first runtime stage. A confident detection is evidence behind a proposal, and evidence grants no permission. Land the segment: a consumer can refuse only what the record lets it check, and it checks before it reads the payload.}
\mlsysfocus{Three invariants of every observation}{%
  {\Large When it was captured, on which clock.\\[0.25cm]
  Where it was measured, in which frame.\\[0.25cm]
  How far it can be trusted, measured.}\\[0.5cm]
  {\normalsize Refuse any claim whose three cannot be checked.}}

% =============================================================================
% CLOSE
% =============================================================================
\begin{frame}{Williston and the rack end, answered}
  \mlsysactivity{Retrieval}{3 min}
  \begin{columns}[c]
    \begin{column}{0.40\textwidth}
      \centering\safeimg[width=\textwidth,height=2.4cm,keepaspectratio]{ntsb-williston.png}\\[0.05cm]
      {\footnotesize Silent omission: no flag to refuse on.}\\[0.2cm]
      {\footnotesize Stopping budget with observation age named: \val{d_rows15} of \val{d_clear_mm} mm}
    \end{column}
    \begin{column}{0.58\textwidth}
      \only<1|handout:0>{\normalsize Name the three things an observation must state, and the two failures today's record cannot catch.}%
      \only<2->{\small \textbf{Rack end:} the frame is \textbf{\val{t_age} ms} old at dispatch, \textbf{\val{d_age} mm} of travel; \val{d_spare15} mm remain for what is not yet priced.\\[0.2cm]
      \textbf{Williston:} every part ran, nothing was detected. A contract exposes faults it can measure; an omission needs independent sensing and coverage.\\[0.2cm]
      \textbf{The record:} capture time and clock, frame, measured error, health.}
    \end{column}
  \end{columns}
  \mlsyscredit{NTSB HAR-17/02. Sensor Perception, \S8.4.1, \S8.7.2, \S8.8; sums: calcs.py.}
  \note{[3 min] [CORE] Retrieval first, 60 seconds of writing. Step 2: back to both hooks. The rack end: an accurate detection is still too late by its age. The worker's position in the frame is \val{t_age} ms old at dispatch, \val{d_age} mm of travel the base has already covered, the fifth row of the stopping budget; with it the partial stop is \val{d_rows15} mm of \val{d_clear} m. The record makes the age checkable; it cannot shorten it. Williston: the system was not designed to, and did not, identify the crossing trailer. No health bit was set because nothing failed. A contract exposes only the faults the ingestion layer measures, so an absent detection is never evidence that the path is clear; the defense is an independent sensing channel with characterized coverage and a conservative permission envelope. The two failures the record cannot catch: a silent semantic omission and a scene outside the support it was qualified on. LINK: Lecture 17 converts detection gaps between groups of people into stopping distance. FLEX: [CORE].}
\end{frame}

\begin{frame}{Takeaways}
  \thumbline{thumb_guard.pdf}{Choose the sensor by the quantity the decision needs; the range that vetoes comes from a sensor the proposer does not use.}\\[0.12cm]
  \thumbline{thumb_errbars.pdf}{Date at capture and place through calibrated frames; a stamp at arrival costs \val{d_arrival} mm, a \val{angle_err}° mount \val{d_deg_far} mm at \val{mount_range} m.}\\[0.12cm]
  \thumbline{thumb_waterfall_step4.pdf}{Observation age, \val{t_age} ms here, is paid in travel, \val{d_age} mm, before anyone reads the frame.}\\[0.12cm]
  \thumbline{thumb_inset_step2.pdf}{The clearance inset rests on measured residuals; learned confidence measures no error in millimeters.}\\[0.12cm]
  \thumbline{thumb_lostmis.pdf}{A misdated frame is worse than a lost one; the contract lets a consumer refuse before reading.}
  \note{[1 min] [CORE] One line per segment. Ask which slide taught each. If short, read the third and the fifth.}
\end{frame}

\begin{frame}{Open questions}
  \begin{columns}[c]
    \begin{column}{0.60\textwidth}
      \small
      \begin{itemize}\setlength{\itemsep}{0.15cm}
        \item End-to-end policies take raw pixels. Should capture age and calibration enter the model as tokens, and how would we test that the model uses them?
        \item How many independent sensing channels does a person at a rack end deserve, and who decides?
        \item A detector qualified on people standing upright may miss a kneeling worker. Whose detection rate does a contract report?
      \end{itemize}
    \end{column}
    \begin{column}{0.38\textwidth}
      \centering\safeimg[width=\textwidth,height=3.4cm,keepaspectratio]{waymo-roof-sensors.jpg}\\[0.05cm]
      {\footnotesize Roof sensor housing on a robotaxi}
    \end{column}
  \end{columns}
  \mlsyscredit{Photo: Waymo Jaguar I-Pace roof sensors, 9yz, CC BY 4.0, Wikimedia Commons. Questions: Sensor Perception, \S8.8, \S8.9, \S8.7.2.}
  \note{[0.5 min] [OPTIONAL] No single answer; pick one for discussion or a short written response. The first comes from the chapter's observation tokenizer and its fallacy about end-to-end policies. Tokens expose time and frame to the model, but only latency-injection ablations show whether it uses them, and the permission path still checks. The second weighs cost against correlated failure. The third previews Lecture 17, which owns detection gaps between groups of people.}
\end{frame}

\begin{frame}{Read and explore}
  \small
  \begin{itemize}\setlength{\itemsep}{0.04cm}
    \item \textbf{Read:} Sensor Perception, \S8.1 to \S8.10; Napkin Math 8.1 to 8.3 by hand; Checkpoints 8.1 to 8.3.
    \item Geiger et al., ``Vision meets Robotics: The KITTI Dataset,'' IJRR 2013: a sensor rig and its timing.
    \item Furgale, Rehder, Siegwart, IROS 2013: estimating camera-to-IMU time offsets.
    \item Hartley and Zisserman, \emph{Multiple View Geometry}, 2003: stereo and its error.
    \item Angelopoulos and Bates, ``Conformal Prediction: A Gentle Introduction,'' 2023.
    \item Gallego et al., ``Event-Based Vision: A Survey,'' TPAMI 2022.
  \end{itemize}
  \note{[0.5 min] [CORE] The chapter first, with its three napkin-math problems done by hand. KITTI's paper documents a real rig: global-shutter cameras, a reed-contact trigger, host-clock stamps. Furgale et al. show why time offsets between sensors must be estimated as unknowns. Hartley and Zisserman derive the stereo geometry. Angelopoulos and Bates explain the conformal rank used for the inset. Gallego et al. survey event cameras. FLEX: an exercise that works well: download one KITTI drive's timestamp files and measure the camera stamp offsets yourself.}
\end{frame}

\begin{frame}[c]
  \begin{columns}[c]
    \begin{column}{0.34\textwidth}
      \centering\safeimg[width=\textwidth,height=4.2cm,keepaspectratio]{thumb_next_stack_brain.pdf}
    \end{column}
    \begin{column}{0.62\textwidth}
      {\small\textcolor{lightgray}{NEXT LECTURE \textbar{} SPATIAL MEMORY}}\\[0.5cm]
      {\LARGE\bfseries How long can a robot act on a scene it can no longer see?}
    \end{column}
  \end{columns}
  \note{[0.5 min] [CORE] Perception answers only for the instant of capture. When the base pivots past a rack upright, the cross aisle leaves the cameras' view, and a worker can step into a corridor the last frame recorded as clear. Lecture 9 fuses observation contracts, each at its own capture time, into a state estimate whose uncertainty grows while no new evidence arrives, and that growth sets the deadline after which a remembered empty aisle can no longer support motion.}
\end{frame}

\end{document}
